% ----- FIG 1 (Cartesian grid) -----

\definecolor{vaxgreen}{RGB}{59,125,35}
\definecolor{contagred}{RGB}{192,0,0}

% Hollow or filled cross node with a number label.
% #1 = coordinate name, #2 = fill colour, #3 = text colour, #4 = label
\newcommand{\crossnode}[4]{%
  \begin{scope}[shift={(#1)}]
    \fill[#2, draw=vaxgreen, line width=1pt]
      (0.15,0.4) -- (0.15,0.15) -- (0.4,0.15) -- (0.4,-0.15) --
      (0.15,-0.15) -- (0.15,-0.4) -- (-0.15,-0.4) -- (-0.15,-0.15) --
      (-0.4,-0.15) -- (-0.4,0.15) -- (-0.15,0.15) -- (-0.15,0.4) -- cycle;
    \node[#3, font=\small\sffamily] at (0,0) {#4};
  \end{scope}
}

% Vaccinated node with a number label
\newcommand{\circnode}[2]{%
  \begin{scope}[shift={(#1)}]
    \fill[vaxgreen!55!white, draw=black, line width=0.8pt] circle (0.42);
    \node[black] at (0,0) {#2};
  \end{scope}
}

% Inoculated node with a number label
\newcommand{\circnodeinoc}[2]{%
  \begin{scope}[shift={(#1)}]
    \fill[vaxgreen!25!white, draw=vaxgreen, line width=1.3pt] circle (0.44);
    \node[black, font=\bfseries] at (0,0) {#2};
  \end{scope}
}

% Directional chevron
% #1 = coordinate, #2 = rotation
\newcommand{\chevron}[2]{%
  \begin{scope}[shift={#1}, rotate=#2]
    \fill[white, draw=vaxgreen, line width=1pt]
      (-0.22,0.28) -- (0.22,0) -- (-0.22,-0.28) -- (-0.08,0) -- cycle;
  \end{scope}
}

% ----- FIG 2 (Strong grid) -----

\definecolor{cYellow}{RGB}{247,230,120}
\definecolor{cRed}{RGB}{224,110,100}
\definecolor{cBlue}{RGB}{120,175,215}
\definecolor{cGreen}{RGB}{110,165,120}

% light wash version for circle fills
\colorlet{fillYellow}{cYellow!55!white}
\colorlet{fillRed}{cRed!60!white}
\colorlet{fillBlue}{cBlue!55!white}
\colorlet{fillGreen}{cGreen!55!white}

% darker text colour for V4, since cYellow is too light to read
\colorlet{cYellowText}{cYellow!70!black}

% Plain contagious circle: #1 = coordinate, #2 = fill colour, #3 = label
\newcommand{\gridcircle}[3]{%
  \begin{scope}[shift={(#1)}]
    \fill[#2, draw=black, line width=0.8pt] circle (0.42);
    \node[black] at (0,0) {#3};
  \end{scope}
}

% Two-tone contagious circle (top/bottom split): #1=coord #2=top fill #3=bottom fill #4=label
\newcommand{\gridsplitcircle}[4]{%
  \begin{scope}[shift={(#1)}]
    \begin{scope}
      \clip (-0.42,0) rectangle (0.42,0.42);
      \fill[#2] (0,0) circle (0.42);
    \end{scope}
    \begin{scope}
      \clip (-0.42,-0.42) rectangle (0.42,0);
      \fill[#3] (0,0) circle (0.42);
    \end{scope}
    \draw[black, line width=0.8pt] (0,0) circle (0.42);
    \node[black] at (0,0) {#4};
  \end{scope}
}

% Vaccinated square (always plain/white): #1 = coordinate, #2 = label
\newcommand{\gridsquare}[2]{%
  \begin{scope}[shift={(#1)}]
    \fill[white, draw=black, line width=0.8pt] (-0.36,-0.36) rectangle (0.36,0.36);
    \node[black] at (0,0) {#2};
  \end{scope}
}

% Root marker
\newcommand{\gridroot}[1]{%
  \begin{scope}[shift={(#1)}]
    \fill[white, draw=black, line width=0.8pt] (-0.36,-0.36) rectangle (0.36,0.36);
    \draw[black, line width=1pt] (-0.22,-0.22) -- (0.22,0.22);
    \draw[black, line width=1pt] (-0.22,0.22) -- (0.22,-0.22);
  \end{scope}
}

% Special inoculation vertex: #1=coord #2=border/fill colour #3=text colour #4=label text (e.g. V_1)
\newcommand{\gridspecial}[4]{%
  \begin{scope}[shift={(#1)}]
    \fill[#2!25!white, draw=#2, line width=1.3pt] circle (0.44);
    \node[#3, font=\bfseries] at (0,0) {$#4$};
  \end{scope}
}

% Chevron (different than Fig. 1 as we need to change colours)
% #1 = coordinate, #2 = rotation (0=right,90=up,180=left,270=down), #3 = colour
\newcommand{\gridchevron}[3]{%
  \begin{scope}[shift={#1}, rotate=#2]
    \fill[white, draw=#3, line width=1pt]
      (-0.22,0.28) -- (0.22,0) -- (-0.22,-0.28) -- (-0.08,0) -- cycle;
  \end{scope}
}

% Bicolour chevron: outline is #3 (colour) on the true top half and #4 (colour) on the true bottom half, regardless of the chevron's own rotation
% #1 = coordinate, #2 = rotation, #3 = top colour, #4 = bottom colour
\newcommand{\gridchevronsplit}[4]{%
  \begin{scope}[shift={#1}]
    \begin{scope}[rotate=#2]
      \fill[white] (-0.22,0.28) -- (0.22,0) -- (-0.22,-0.28) -- (-0.08,0) -- cycle;
    \end{scope}
    \begin{scope}
      \clip (-0.5,0) rectangle (0.5,0.5);
      \begin{scope}[rotate=#2]
        \draw[#3, line width=1pt] (-0.22,0.28) -- (0.22,0) -- (-0.22,-0.28) -- (-0.08,0) -- cycle;
      \end{scope}
    \end{scope}
    \begin{scope}
      \clip (-0.5,-0.5) rectangle (0.5,0);
      \begin{scope}[rotate=#2]
        \draw[#4, line width=1pt] (-0.22,0.28) -- (0.22,0) -- (-0.22,-0.28) -- (-0.08,0) -- cycle;
      \end{scope}
    \end{scope}
  \end{scope}
}

% map symbolic names to drawing commands
\newcommand{\cellyellow}[2]{\gridcircle{#1}{fillYellow}{#2}}
\newcommand{\cellred}[2]{\gridcircle{#1}{fillRed}{#2}}
\newcommand{\cellblue}[2]{\gridcircle{#1}{fillBlue}{#2}}
\newcommand{\cellgreen}[2]{\gridcircle{#1}{fillGreen}{#2}}

\newcommand{\cellS}[2]{\gridsquare{#1}{#2}}
\newcommand{\cellX}[2]{\gridroot{#1}}

\newcommand{\cellVone}[2]{\gridspecial{#1}{cRed}{black}{V_1}}
\newcommand{\cellVtwo}[2]{\gridspecial{#1}{cBlue}{black}{V_2}}
\newcommand{\cellVthree}[2]{\gridspecial{#1}{cGreen}{black}{V_3}}
\newcommand{\cellVfour}[2]{\gridspecial{#1}{cYellow}{black}{V_4}}